\honest{Dreams of AI Design}{Eliezer Yudkowsky}{2008}

Quite a few ``AGI wannabes,'' I say, seem to think that living inside a mind for a decade or two teaches you how minds work. It does not. ``Artificial Intelligence is fundamentally about reducing the mental to the non-mental.'' \nb{Dennett had put the same idea in 1978: progress is replacing an intelligent homunculus with a team of ``relatively ignorant, narrow-minded, blind'' ones.} Introspection cannot teach this, because the brain's work is hidden and nothing tells you it has been handed off.

Did Aristotle realize that his final causes asked, ``What would this object do, if it could make plans?'' ``I rather doubt it.'' \nb{Aristotle did ask whether natural purposes need a planner, and answered no: ``Art does not deliberate.''} Aristotle thought the brain cooled the blood, and that humans, with larger brains, were calmer and more contemplative. So cool the computer! That explains nothing. The only way to predict what a contemplative thing does, I say, is to imagine being one: empathic inference. \nb{As in ``Humans in Funny Suits,'' this is simulation theory stated as fact; its rival holds that people use a tacit theory of other minds.} Saying that many neurons produce contemplation by ``emergence'' is no better. It draws an arrow between two labelled boxes.

Real reductions look different. Jaynes takes you from the feeling of evidence to Bayes's theorem. Pearl explains why, after your burglar alarm goes off, news of an earthquake makes you doubt the burglary. \nb{The alarm is Pearl's own example.} The minimax search tree tells a chess program what a good move is. Until you know these, you may think good moves come from cooling the blood.

Many people I know believe intelligence is ``the product of'' common-sense knowledge, massive parallelism, creative destruction or intuition. They build detached levers, LISP tokens labelled apple, or ``massively parallel'' neural nets ``just like the human brain,'' and are shocked when nothing much happens. \nb{The post names none of them. Four years later a deep neural network trained on graphics processors won a major image-recognition contest, and deep learning followed; the post had mocked nets justified by likeness to the brain, not neural networks as such.} Designs drawn as boxes with meaningful labels are dreams; for ``a truly epic example,'' see any Mentifex diagram. \nb{Mentifex is Arthur T. Murray, author of \textsc{AI4U}.}

My rule: if someone says their AI will do something ``just like humans,'' give them no venture capital until they explain what it does and why without pointing at humans. Planes are ``not kept aloft by analogies.'' \nb{Robin Hanson objected in the comments that the rule would refuse whole brain emulation, which copies a brain rather than arguing from analogy. The author replied that it ``might be physically possible'' but was not worth funding.}

This art is ``demonstrably'' difficult, I say, and I give no demonstration. That is why I have written at such length about reductionism, Taboo and anthropomorphism.
